\chapter{Ergebnisse und Diskussion}
\label{chap:ergebnisse_und_diskussion}

\section{Ergebnisse des Validierungssplits}

Die Ergebnisse des Validierungssplits werden im späteren Verlauf für die Auswahl der besten Modellkonfiguration verwendet.

Tabelle \ref{tab:validation_ergebnisse} fasst die wichtigsten Rankingmetriken des Referenzlaufs A08 auf dem Validierungssplit zusammen. 

\begin{table}[tbp] 
  \centering 
  \caption{Ergebnisse des Referenzlaufs A08 auf dem Validierungssplit} 
  \label{tab:validation_ergebnisse} 
  \begin{tabular}{lr} \toprule \textbf{Metrik} & \textbf{Wert} \\ 
    \midrule 
    Target Recall@1 & 0,3137 \\
    Target Recall@3 & 0,4472 \\
    Target Recall@5 & 0,5031 \\
    Target Recall@10 & 0,6211 \\
    Family Recall@1 & 0,8199 \\
    Family Recall@10 & 0,8913 \\
    MRR & 0,4121 \\
    Hierarchical nDCG@10 & 0,6766 \\
    \bottomrule 
  \end{tabular} 
\end{table}

Die Ergebnisse zeigen, dass das Modell bereits nach dem Test durch den Validierungssplit die korrekte Makrofamilie besser identifiziert als die exakte Ziel-Makro \parencite{EPLAN-Makros}.
Während das korrekte Ziel-Makro \parencite{EPLAN-Makros} in 31,37\,\% der Anfragen auf Platz 1 steht, wird die korrekte Makrofamilie in 81,99\,\% der Anfragen auf die erste Position gesetzt.

Durch den Target Recall@10 lässt sich erkennen, dass das Modell in 62,11\,\% der Anfragen das korrekte Ziel-Makro \parencite{EPLAN-Makros} innerhalb der Top-10-Vorschläge platziert.
Interessant ist der Abstand zwischen Family Recall@10 (89,13\,\%) und Target Recall@10 (62,11\,\%). Dies deutet darauf hin, dass viele Fehlkategorisierungen innerhalb derselben Makrofamilie stattfinden.

Die erzielten Ergebnisse bei MRR von 0,4121\,\% sowie ein hierarchisches nDCG@10 von 0,6766\,\%  zeigen zusätzlich, dass beim Validierungssplit das Modell in der Lage ist eine sinnvolle Reihenfolge zu erzeugen, 
es jedoch weiterhin Schwierigkeiten bei der korrekten Auflösung technischer Varianten besitzt.\parencite{LTR_Configuration_Test_Series}

\section{Ergebnisse des Testsplits}

Nach dem die Modellauswahl abgeschlossen wurde, erfolgt die finale Evaluation des ausgewählten Modells auf dem unberührten Testsplit.

Tabelle \ref{tab:test_ergebnisse} zeigt die erzielten Ergebnisse des Referenzlaufs A08 des Testsplits. 

\begin{table}[tbp] 
  \centering 
  \caption{Ergebnisse des Referenzlaufs A08 auf dem Testsplit} 
  \label{tab:test_ergebnisse} 
  \begin{tabular}{lr} \toprule \textbf{Metrik} & \textbf{Wert} \\ \midrule 
    Target Recall@1 & 0,3478 \\
    Target Recall@3 & 0,4783 \\
    Target Recall@5 & 0,5494 \\
    Target Recall@10 & 0,6798 \\
    Family Recall@1 & 0,8735 \\
    Family Recall@10 & 0,9368 \\
    MRR & 0,4534 \\
    Hierarchical nDCG@10 & 0,7449 \\
    \bottomrule 
  \end{tabular} 
\end{table}

Beim Testsplit zeigt sich ein ähnliches Bild wie beim Validierungssplit. Das Modell sagt bei 34,78\,\% der Anfragen das korrekte Ziel-Makro \parencite{EPLAN-Makros} auf Platz 1 voraus, während die korrekte Makrofamilie in 87,35\,\% der Anfragen auf die erste Position gesetzt wird.

Wie beim Validierungsplit ist auch hier der deutliche Abstand zwischen Family Recall@10 (93,68\,\%) und Target Recall@10 (67,98\,\%) zu erkennen. 

Der erreichte MRR- Wert von 45,34\,\% sowie die hierarchische nDCG@10 von 74,49\,\% bestätigen zusätzlich die Erkenntnisse aus dem Validierungssplit.\parencite{LTR_Configuration_Test_Series}

\section{Fehleranalyse}

Um die Modellleistung detailliert erfassen zu können, werden die Top-1-Ergebnisse des Testsplits analysiert. Dabei werden die fachlichen Beziehungen zum tatsächlichen Zielmakro untersucht.

Tabelle \ref{tab:fehleranalyse_top1} zeigt die Verteilung der einzelnen Fehlerkategorien. 

\begin{table}[tbp] 
  \centering 
  \caption{Top-1-Analyse der Vorhersagen auf dem Testdatensatz} 
  \label{tab:fehleranalyse_top1} 
  \begin{tabular}{lrr} 
    \toprule 
    \textbf{Kategorie} & \textbf{Anzahl} & \textbf{Anteil} \\ 
    \midrule 
    Exaktes Ziel-Makro \parencite{EPLAN-Makros} & 88 & 34,78\,\% \\
    Gleiche Familie, falsches Target & 133 & 52,57\,\% \\
    Gleicher Familienpräfix & 12 & 4,74\,\% \\
    Andere Familie & 20 & 7,91\,\% \\
    \bottomrule 
  \end{tabular} 
\end{table}

Die Fehleranalyse bestätigt die Beobachtungen aus den Metriken. Von insgesamt 253 Testanfragen werden 133 Fälle einer gleichen Familie, aber einer falschen Variante zugeordnet, was 52,57\,\% entspricht.
In 7,91\,\% der Fälle wird eine andere Familie vorhergesagt. Das zeigt, dass das Modell die grundlegenden technischen Zusammenhänge zwischen Artikeln und Makrofamilien zuverlässig erlernt hat.
Die Analyse zeigt, dass die zentrale Herausforderung die Identifikation des korrekten Ziel-Makros \parencite{EPLAN-Makros} innerhalb der richtigen Makrofamilie ist. Dafür muss die Variantenauflösung innerhalb der Makrofamilie weiter verbessert werden.
Die große Differenz zwischen dem Target Recall@1 und dem Family Recall@1 unterstreicht diese Herausforderung zusätzlich.
Dies wird im Kapitel \ref{chap:grenzen_und_ausblick} genauer analysiert.\parencite{LTR_Configuration_Test_Series}


\section{Rangstatistische Analyse}

Eine weitere wichtige Analyse ist die statistische Auswertung der erzeugten Ranglisten.


Die wichtigsten Rangstatistiken des Testdatensatzes werden in Tabelle \ref{tab:rangstatistik} dargestellt.

\begin{table}[tbp] 
  \centering 
  \caption{Rangstatistische Kennzahlen des Testsplits} 
  \label{tab:rangstatistik} 
  \begin{tabular}{lr} 
    \toprule 
    \textbf{Kennzahl} & \textbf{Wert} \\ 
    \midrule 
    Erwarteter Zielrang (Mean Rank) & 23,96 \\
    Median Rank & 4 \\
    90\,\%-Perzentil Rank & 61,6 \\
    Maximaler Rang & 360 \\
    \bottomrule 
  \end{tabular} 
\end{table}

Der Median Rank von 4 zeigt, dass das korrekte Makro \parencite{EPLAN-Makros} in der Hälfte der Testanfragen unter den ersten vier Positionen der Rangliste erscheint.
Der mittlere Rang (Mean Rank) von 23,96 zeigt gleichzeitig, dass es einzelne schwierige Anfrage gibt, die deutliche schlechter eingeordnet werden.
Diese Annahme wird durch das 90\,\%-Perzentil von 61,6 bestätigt. Dabei sitzt das Makro \parencite{EPLAN-Makros} in 90\,\% der Testanfragen unter den ersten 61,6 Positionen der Rangliste.

Auch die Rangstatistiken zeigen, dass das Modell in der Lage ist, die korrekte Makro-Familie in den meisten Fällen weit oben in der Rangliste zu platzieren.\parencite{LTR_Configuration_Test_Series}

\section{Analyse der Generalisierungsfähigkeit}

Zur Beurteilung der Generalisierungsfähigkeit des Modells wurde eine zusätzliche Leave-One-Out Kontrollanalyse durchgeführt. 
Um eine Verfälschung der Ergebnisse der Trainingsanalyse zu verhindern, werden bei einer Anfrage genau dieser Artikel aus der Anfragemenge ausgeschlossen.


Tabelle \ref{tab:loo_vergleich} stellt die Ergebnisse der Leave-One-Out-Analyse (LOO) sowie der Validierungs- und Testergebnisse gegenüber.

\begin{table}[tbp] 
  \centering 
  \caption{Vergleich von regulärer Trainingsbewertung, Leave-One-Out-Analyse, Validation und Test} 
  \label{tab:loo_vergleich} 
  \begin{tabular}{lrrrr} 
    \toprule 
    \textbf{Metrik} & \textbf{Train (LOO)} & \textbf{Validation} & \textbf{Test} \\ 
    \midrule 
    Recall@1 & 0,1870 & 0,3137 & 0,3478 \\
    Recall@3 & 0,3028 & 0,4472 & 0,4783 \\
    Recall@10 & 0,4695 & 0,6211 & 0,6798 \\
    MRR & 0,2864 & 0,4121 & 0,4534 \\
    Family Recall@1 & 0,7246 & 0,8199 & 0,8735 \\
    Hierarchical nDCG@10 & 0,5945 & 0,6766 & 0,7449 \\
    Mean Rank & 25,93 & 23,96 & 23,96 \\
    Median Rank & 13,0 & 4,0 & 4,0 \\
    90\,\%-Perzentil & 62,0 & 61,6 & 61,6 \\
        \bottomrule 
    \end{tabular} 
    \end{table}

Es zeigt sich, dass die Leave-One-Out-Bewertung eine strengere Evaluationsmethode als die Bewertung durch den Validierungs- und Testdatensatz darstellt.
Die Ursache dafür liegt im Support der Kandidatenprofile. Von den vorhanden 375 Kandidatenprofilen basieren insgesamt 302 Profile ausschließlich auf einem Trainingsartikel.
Wenn dieser Artikel aus der Anfragemenge ausgeschlossen wird, enthält das entsprechende Kandidatenprofil keine unterstützenden Trainingsdaten mehr und kann dadurch keinen Beitrag zur korrekten Vorhersage leisten.

Bei den Validierungs- und Testdaten tritt dieser Effekt nicht auf, da durch die im Datensplit erzwungene Trainingsabdeckung jedes Ziel-Makro \parencite{EPLAN-Makros} mindestens einen Artikel in jedem Split enthält, welcher nicht identisch zu den andern beiden ist.


\section{Vergleich der Konfigurationsvarianten}

Um verschiedene Modellkonfigurationen zu vergleichen, wurde die aktuell Versuchsserie A01 bis A09 durchgeführt. 
Dabei wurden unterschiedliche Strategien, Profilbreiten und Regularisierungsparameter getestet. Ziel davon ist es, die Einflüsse einzelner Parameter auf die Modellleistung zu analysieren.\parencite{LTR_Configuration_Test_Series}

\subsection{Untersuchte Konfigurationsparameter}

Die verschiedenen Varianten der Konfiguration haben unterschiedliche Auswirkungen auf die Trainingslabels, die Kandidatenprofile sowie der Modellregularisierung. 
Die untersuchten Konfigurationsparameter werden in folgender Tabelle \ref{tab:konfigurationsparameter} erläutert.

\begin{table}[tbp] 
    \centering 
    \caption{Untersuchte Konfigurationsparameter} 
    \label{tab:konfigurationsparameter} 
    \begin{tabularx}{\textwidth}{p{4cm}X} 
        \toprule \textbf{Parameter} & \textbf{Bedeutung} \\ 
        Same-Family-Gewicht &  Negative Kandidaten aus derselben Makrofamilie werden innerhalb der Verlustfunktion unterschiedlich gewichtet. Dadurch kann die Variantenauflösung stärker betont werden. \\ 
        Profilbreite & Begrenzt die Anzahl der gespeicherten Profilwerte innerhalb eines Kandidatenprofils. Kleinere Profilbreiten erzeugen eine kompaktere Kandidatenrepräsentation. \\ 
        L1-Regularisierung ($\alpha$) & Bestraft große Modellgewichte und reduziert die Modellkomplexität. Dadurch werden irrelevante Merkmale stärker unterdrückt. \\ 
        L2-Regularisierung ($\lambda$) & Stabilisiert die Gewichtung der einzelnen Merkmale und reduziert die Gefahr einer Überanpassung an einzelne Trainingsbeispiele. \\ 
        Min. Split Gain & Definiert den minimal notwendigen Informationsgewinn um eine Baumaufteilung im LightGBM-Modell zu ermöglichen. Höhere Werte führen dadurch zu konservativeren Entscheidungsbäumen. \\ 
        \bottomrule 
    \end{tabularx} 
\end{table} 

A02 bis A07 wurden so gewählt, dass immer nur ein Parameter gezielt verändert wurde. Dadurch lassen sich die Auswirkungen der einzelnen Konfigurationsparameter isoliert untersuche.
Bei A08 und A09 werden mehrere Parameter (L1 und L2) gleichzeitig verändert, da diese bei LightGBM konzeptionell derselbe Hebel sind (Bestrafung der Blattgewichte). 
Um eine sinnvolle "stärkere Regularisierung" zu erzielen, ist es in der Praxis sinnvoll, beide gemeinsam hochzudrehen, nicht nur einen isoliert.
Bei A09 kommt zusätzlich ein positiver Wert für \texttt{min\_split\_gain} zum Einsatz um dadurch einen noch größeren Einfluss auf eine "stärkere Regularisierung" zu erzielen.\parencite{LTR_Configuration_Test_Series}


\subsection{Referenzkonfiguration A01}

A01 dient als Referenzkonfiguration und bildet die Basis für die Vergleichsexperimente. Sie enthält die Standardwerte für alle untersuchten Konfigurationsparameter.
Alle weiteren Konfigurationen A02 bis A09 werden gegenüber dieser bewertet.

Die Referenzparameter lauten: 

\begin{table}[tbp] 
    \centering 
    \caption{Referenzkonfiguration A01} 
    \label{tab:a01_referenz} 
    \begin{tabular}{ll} 
        \toprule 
        \textbf{Parameter} & \textbf{Wert} \\ 
        \midrule 
        Exaktes Targettraining & Ja \\ 
        Same-Family-Gewicht & 5,0 \\ 
        Profilbreite & 16 \\ 
        L1-Regularisierung ($\alpha$) & 0,05 \\ 
        L2-Regularisierung ($\lambda$) & 0,2 \\ 
        Min. Split Gain & 0,0 \\ 
        \bottomrule 
    \end{tabular} 
\end{table}

Ausgehen von dieser Basis werden die folgenden Varianten der Konfigurationen A02 bis A09 verändert.

Tabelle \ref{tab:konfigurationsabweichungen} zeigt die gegenüber A01 vorgenommenen Veränderungen. 
\begin{table}[tbp] 
    \centering 
    \caption{Abweichungen der Konfigurationen gegenüber der Referenz A01} 
    \label{tab:konfigurationsabweichungen} 
    \begin{tabular}{lp{8cm}} 
        \toprule 
        \textbf{Konfiguration} & \textbf{Änderung gegenüber A01} \\ 
        \midrule 
        A02 & Same-Family-Gewicht von 5,0 auf 2,0 reduziert \\ 
        A03 & Same-Family-Gewicht von 5,0 auf 8,0 erhöht \\ 
        A04 & Same-Family-Gewicht von 5,0 auf 12,0 erhöht \\ 
        A05 & Hierarchische Relevanzlabels statt ausschließlich exakter Zielmakros \\ 
        A06 & Profilbreite von 16 auf 12 reduziert \\ 
        A07 & Profilbreite von 16 auf 8 reduziert \\ 
        A08 & L1-Regularisierung auf 0,1 und L2-Regularisierung auf 0,5 erhöht \\ 
        A09 & L1-Regularisierung auf 0,5, L2-Regularisierung auf 1,0 und \texttt{min\_split\_gain} auf 0,01 erhöht \\ 
        \bottomrule 
    \end{tabular} 
\end{table}

Dadurch lässt sich der Einfluss der jeweiligen Parameteränderung der entsprechenden Konfiguration zuordnen.\parencite{LTR_Configuration_Test_Series}

\subsection{Ergebnisse der Vergleichskonfigurationen}

Tabelle \ref{tab:konfigurationsvergleich} zeigt die getesteten Modellkonfigurationen und deren Leistung im Vergleich.

\begin{table}[tbp] 
    \centering 
    \caption{Ergebnisse der A01--A09-Konfigurationen auf dem Testdatensatz} 
    \label{tab:konfigurationsvergleich} 
    \small 
    \begin{tabular}{lcccc} 
        \toprule 
        \textbf{Konfiguration} & \textbf{Recall@1} & \textbf{Recall@3} & \textbf{Recall@10} & \textbf{MRR} \\ 
        \midrule 
        A01 Referenz & 0,3241 & 0,4822 & 0,6798 & 0,4418 \\ 
        A02 Gewicht 2 & 0,3162 & 0,4862 & 0,6680 & 0,4399 \\ 
        A03 Gewicht 8 & 0,3241 & 0,4743 & 0,6838 & 0,4397 \\ 
        A04 Gewicht 12 & 0,2925 & 0,4783 & 0,6759 & 0,4195 \\ 
        A05 Hierarchisch & 0,3162 & 0,4783 & 0,6719 & 0,4328 \\ 
        A06 Profilbreite 12 & 0,3162 & 0,4862 & 0,6877 & 0,4374 \\ 
        A07 Profilbreite 8 & 0,3241 & 0,4862 & 0,6719 & 0,4439 \\ 
        A08 Regularisierung mittel & \textbf{0,3478} & 0,4783 & 0,6798 & \textbf{0,4534} \\ 
        A09 Regularisierung stark & 0,3399 & 0,4783 & \textbf{0,6917} & 0,4505 \\ 
        \bottomrule 
    \end{tabular} 
\end{table}


\subsection{Einfluss des Same-Family-Gewichts}

Mit den Konfigurationen A02 bis A04 wird der Einfluss des Same-Family-Gewichts untersucht. 
Die Ergebnisse zeigen, dass eine Verringerung des Gewichts auf 2 den Recall@1 auf 31,62\,\% senkt, während eine Erhöhung auf 8 keine Auswirkung zeigt.
Die Gewichtung auf 12 führt zu einer deutlichen Verschlechterung des Recall@1 auf 29,25\,\%.

Daraus lässt sich schließen, dass ein zu hohes Same-Family-Gewicht die Lernstabilität negativ beeinflusst.

\subsection{Einfluss der hierarchischen Labeldefinition}

A05 untersucht die hierarchische Relevanzdefinition. Makros \parencite{EPLAN-Makros} derselben Familie und desselben Präfixes werden dabei hierarchisch relevant betrachtet.
Es zeigt sich jedoch keine Verbesserung gegen dem exakten Targettraining. Der Recall@1 sinkt von 32,41 \,\% auf 31,62 \,\% und die MRR von 44,18 \,\% auf 43,28 \,\%.
Die dadurch entstandene reichhaltigere Zieldefinition hat somit keinen positiven Einfluss auf die Modellleistung.\parencite{LTR_Configuration_Test_Series}

\subsection{Einfluss der Profilbreite}

In den Konfigurationen A06 und A07 werden unterschiedliche Profilbreiten untersucht. Bei A06 beträgt diese zwölf bei A07 acht.
Die Erhöhung der Profilbreite führt zu einer Verschlechterung des Recall@1 von 32,41 \,\% auf 31,62 \,\% und der MRR von 44,18 \,\% auf 43,74 \,\%.
Eine Verringerung der Profilbreite führt zu keiner nennenswerten Veränderung der Modellleistung. 
Die Ergebnisse zeigen, dass sich eine verringerte Profilbreite die Modellleistung nicht signifikant verbessert.\parencite{LTR_Configuration_Test_Series}

\subsection{Einfluss der Regularisierung}

Die deutlichsten Verbesserungen werden durch die Veränderung der Regularisierung erzielt.
A08 verwendet eine mittlere L1- und L2-Regularisierung (\(\alpha = 0{,}1\), \(\lambda = 0{,}5\)). 
Dadurch wird mit einem Recall@1 von 0,3478 der höchste Top-1-Wert aller untersuchten Konfigurationen erreicht. Auch der höchste MRR Wert von 0,4534 wird durch A08 erreicht.
Mit A09 wird die Regularisierung nochmals deutlich erhöht (\(\alpha = 0{,}5\), \(\lambda = 1{,}0\)). Zusätzlich wird ein positiver Wert für \texttt{min\_split\_gain} gesetzt, um die Stabilität der Baumaufteilung zu erhöhen. 
Diese Konfiguration erreicht den höchsten Recall@10 von 0,6917 und die zweithöchste MRR von 0,4505.

Die L1/L2-Regularisierung wirkt sich bei LightGBM direkt auf die Blattgewichte aus.
Sie dämpft diese bei geringer Aussagekraft auf null und unterdrückt dadurch das Rauschen der 2.381 Merkmale bei nur 984 Trainingsbeispielen, ohne das schwache Unterscheidungssignal zwischen Makrofamilien-Varianten zu verlieren.
Bei A09 wird durch die stärkere Dämpfung auch dieses Signal mit reduziert, was die Top-1-Präzision verringert, dabei aber die Gesamtrangfolge robuster macht.\parencite{LTR_Configuration_Test_Series}

\subsection{Beste Konfiguration der Studie}

Die Beste Konfiguration der Studie ist A08, da sie sowohl den höchsten Recall@1 als auch den höchsten MRR erreicht.\parencite{LTR_Configuration_Test_Series}

Die wichtigsten Kennzahlen dieser Variante sind in Tabelle \ref{tab:a08_best} dargestellt. 

\begin{table}[tbp] 
    \centering 
    \caption{Ergebnisse der besten Konfiguration A08} 
    \label{tab:a08_best} 
    \begin{tabular}{lr} 
        \toprule 
        \textbf{Kennzahl} & \textbf{Wert} \\ 
        \midrule Recall@1 & 0,3478 \\ 
        Family Recall@1 & 0,8735 \\ 
        Recall@10 & 0,6798 \\ 
        Family Recall@10 & 0,9368 \\
        MRR & 0,4534 \\ 
        Mittlerer Zielrang & 23,96 \\ 
        90\,\%-Perzentil & 61,6 \\ 
        Schlechtester Rang & 360 \\ 
        \bottomrule 
    \end{tabular} 
\end{table}

Die Fehleranalyse dieser Konfiguration zeigt, dass 7,9\,\% aller Anfragen einer fachlich falschen Makro-Familie zugeordnet wurden.
52,6\,\% der Fehlklassifikationen erfolgen innerhalb der korrekten Makrofamilie. 
Die größte Herausforderung liegt damit weiterhin in der Variantenauflösung technisch ähnlicher Makrofamilien.\parencite{LTR_Configuration_Test_Series}

